% Stile del report: tutte le scelte visive sono raccolte qui.
% I capitoli contengono contenuti, tabelle e immagini, senza temi separati.
\usepackage{fontspec}
\defaultfontfeatures{Ligatures=TeX,Scale=MatchLowercase}
\usepackage[italian]{babel}
\usepackage[a4paper,margin=2.5cm]{geometry}
\usepackage{xcolor}
\usepackage[breakable]{tcolorbox}
\usepackage{amsmath,amssymb}
\usepackage{microtype}
\usepackage{longtable,booktabs,array}
% Compatibilita con le tabelle storiche delle fasi C/D, non riscritte qui.
\usepackage{calc}
\newcounter{none}
\usepackage{graphicx}
\usepackage{pgfplots}
\pgfplotsset{compat=1.18}
\usepackage{caption}
\usepackage{enumitem}
\usepackage{needspace}
\usepackage{hyperref}
\usepackage{bookmark}
\usepackage{xurl}
\usepackage{fancyhdr}
\usepackage{titlesec}
\usepackage{csquotes}
\usepackage[backend=biber,style=numeric-comp,sorting=nyt,maxbibnames=99,defernumbers=true,urldate=long]{biblatex}
\DefineBibliographyStrings{italian}{urlseen={consultato il}}
\addbibresource{references/references.bib}

% Testo, elenchi e tabelle: poche regole comuni, senza spaziature per pagina.
\setcounter{secnumdepth}{2}
\setcounter{tocdepth}{2}
\setlength{\emergencystretch}{6em}
\setlength{\parindent}{0pt}
\setlength{\parskip}{0.4\baselineskip plus 1pt minus 1pt}
\setlength{\LTpre}{0.6\baselineskip}
\setlength{\LTpost}{0.6\baselineskip}
\renewcommand{\arraystretch}{1.15}
\setlength{\tabcolsep}{5pt}
\setlist[itemize]{leftmargin=*,topsep=0.4\baselineskip,itemsep=0.15\baselineskip,parsep=0pt,partopsep=0pt}
\captionsetup{font=small,labelfont=bf,skip=6pt}
\urlstyle{same}
\AtBeginDocument{\renewcommand{\abstractname}{Abstract}}
\raggedbottom
% Evita righe isolate all'inizio e alla fine dei paragrafi.
\clubpenalty=10000
\widowpenalty=10000
\displaywidowpenalty=10000

% Un solo colore di accento; i contenuti restano neri su bianco.
\definecolor{headingcolor}{HTML}{24364B}
\definecolor{rulecolor}{HTML}{A7B4C2}
\definecolor{todobackground}{HTML}{FFF4CC}
\definecolor{todotext}{HTML}{6B3E00}
\titleformat{\chapter}{\normalfont\huge\bfseries\raggedright\color{headingcolor}}{\thechapter}{0.6em}{}
\titlespacing*{\chapter}{0pt}{0pt}{1.5\baselineskip}
\titleformat{\section}{\normalfont\Large\bfseries\raggedright\color{headingcolor}}{\thesection}{0.6em}{}
\titlespacing*{\section}{0pt}{1.6\baselineskip}{0.7\baselineskip}
\titleformat{\subsection}{\normalfont\large\bfseries\color{headingcolor}}{\thesubsection}{0.6em}{}
\titlespacing*{\subsection}{0pt}{\baselineskip}{0.4\baselineskip}
\titleformat{\subsubsection}{\normalfont\normalsize\bfseries}{\thesubsubsection}{0.6em}{}
\titlespacing*{\subsubsection}{0pt}{0.8\baselineskip}{0.3\baselineskip}

% Le card scorrono come sezioni; i campi restano uniti al testo successivo.
\newcommand{\cardsection}[1]{\par\Needspace{11\baselineskip}\section{#1}}
\newcommand{\fieldheading}[1]{%
  \par\Needspace{5\baselineskip}%
  \par\addvspace{0.65\baselineskip}%
  {\noindent\color{headingcolor}\bfseries #1\par}%
  \nobreak\vspace{0.1\baselineskip}%
}
\newtcolorbox{personacard}[1]{
  breakable,colback=white,colframe=rulecolor,
  colbacktitle=white,coltitle=headingcolor,fonttitle=\bfseries,
  title={#1},boxrule=0.5pt,arc=0pt,
  left=7pt,right=7pt,top=5pt,bottom=5pt,
  before skip=0.5\baselineskip,after skip=0.5\baselineskip
}

% Annotazioni dei contenuti ancora da completare nelle fasi C/D.
\newtcolorbox{projecttodobox}{
  breakable,colback=todobackground,colframe=todotext,
  colbacktitle=todobackground,coltitle=todotext,
  title={TODO},fonttitle=\bfseries,boxrule=0.5pt,arc=0pt,
  left=7pt,right=7pt,top=5pt,bottom=5pt
}
\newcommand{\projecttodo}[1]{\begin{projecttodobox}#1\end{projecttodobox}}
\newcommand{\todoinline}[1]{\textcolor{todotext}{\textbf{TODO:} #1}}

\hypersetup{colorlinks=true,linkcolor=headingcolor,citecolor=headingcolor,
  urlcolor=headingcolor,pdftitle={Documentazione del progetto HPHP},pdfauthor={Team Affordance}}
\pagestyle{fancy}
\fancyhf{}
\fancyhead[L]{\small\color{headingcolor} HPHP / CUPweb}
\fancyfoot[C]{\thepage}
\setlength{\headheight}{14pt}
\renewcommand{\headrulewidth}{0pt}
\fancypagestyle{plain}{\fancyhf{}\fancyfoot[C]{\thepage}\renewcommand{\headrulewidth}{0pt}}
